% =====================================================================================
% sections/s8_discussion.tex -- Discussion, limitations and conclusion.
% No computation and no new numbers: every figure quoted here is stated or tabulated in
% an earlier section.  The "% src:" comments name the underlying file (paths relative to
% the root of the repository).
% =====================================================================================
\section{Discussion, limitations and conclusion}\label{s8_discussion:sec}

\subsection{What the evidence supports}\label{s8_discussion:sec:evidence}

Findings (1) to (5) are ordered by the strength of their evidence, strongest first; (6) and (7)
concern Sections~\ref{s6_highdim:sec:bound} and~\ref{s6_highdim:sec:remedies}, and (8) the
reporting of errors.

\begin{enumerate}[label=(\arabic*),leftmargin=*]
% src: data/s7_pitfalls_autograd.json (operator_unit_test_float64; training.max_pde_term_over_all_steps 0.0); research_benchmark/results/oneface_laplace3d.json
\item \emph{Two implementation faults give a small loss and a wrong answer.} A residual made
identically zero by a misused automatic-differentiation call, and a boundary-value problem with data
on one face of six, each gave a small training loss and a wrong answer (Section~\ref{s7_pitfalls:sec}): in the
first the network follows whatever the data term prescribes, which a correct residual would have
penalised. Neither fault depends on seeds or budgets, and each is exposed by a cheap test.
% src: research_benchmark/results/tables.md (T2: 175 of 175); data/s4_lowdim_numbers.json (lbfgs_vs_logged_adam_iterates_it1500_3000); verify_research_benchmark/v_summary.out (39 of 40); data/s5_hard_wave2d_long.json (summary.per_seed: adam_lbfgs_rel_l2 5.7e-3 to 9.3e-3)
\item \emph{An \LBFGS\ phase is the most effective change tested.} At equal iteration count it
lowered the held-out error in all 175 pairs (39 of 40 in the verification code), by about one order
of magnitude on \PH\ and \PLtwo\ even against the best logged \Adam\ iterate, and a longer phase took
\PWtwo\ below one per cent (Sections~\ref{s4_lowdim:sec:optimiser} and~\ref{s5_hard:sec:w2}). The
\Adam\ arm was not tuned.
% src: research_benchmark/results/tables.md (T3, T6; T3 laplace2d adam resample 6.36, 0/10, the largest ratio of the table); data/s4_lowdim_gridcontrol.json (column against placebo 0.19, 10 of 10); verify_research_benchmark/v_summary.out
\item \emph{The choice among space-filling point sets is secondary.} No strategy is best on every
problem; the largest effect after the \LBFGS\ phase, a fourfold penalty of the cell-centred grid on
\PH, is caused by the absence of residual points next to the initial line
(Table~\ref{s4_lowdim:tab:gridcontrol}). Adaptive sampling brought no detectable gain at our budget,
and redrawing the points at every iteration hurt under \Adam\ alone (on \PLtwo\ by a factor 6.4,
the largest penalty of any strategy in the benchmark) but not after the \LBFGS\ phase (Section~\ref{s4_lowdim:sec:strategy}).
% src: research_benchmark/results/benchmark_runs.csv (wave1d: 50 runs, 0.414 to 0.512); data/s5_hard_wave1d_long.json (summary.per_seed: 0.337 to 0.381); data/s5_hard_cost.json (summary: ratio_smallest_of_all_measurements 475.6, ratio_B_range 1.9e3 to 1.3e5)
\item \emph{At the sizes and budgets used, the unit-weight PINN does not solve the two-mode wave
problem, and finite differences are far faster.} All 50 runs on \PWone\ ended between 41 and 51 per
cent error and the longer runs between 34 and 38 per cent, although \citet{rathore2024} report a few
per cent with wider networks, more points and a tuned schedule. A second-order finite-difference solve
was at least as accurate as the best PINN on all five problems and faster by a factor of more than
400 in every timing (smallest measured ratio $\sci{4.8}{2}$; Section~\ref{s5_hard:sec:fd}), as \citet{grossmann2024} found against finite
elements.
% src: research_highdim/results/summary.csv; research_highdim/results/extensions.json; data/s6_highdim_long_d20.jsonl; data/s6_highdim_robin_bias.json (against_trained_networks.beta_100_16000_iterations: bias_over_gap 1.048, 0.654); data/s6_highdim_bound_numbers.json (recheck_r2: P1_ritz_eqcpu_d20, P1_pinn_d20)
\item \emph{In higher dimension the ranking of PINN and Deep Ritz depends on what is counted.} Per
iteration the PINN was the more accurate for $d\le10$ and, with 16\,000 iterations, at $d=20$; at
equal compute Deep Ritz was the more accurate (at $d=10$, and at $d=20$ for \Pone\ in a descriptive
re-check with three seeds; Sections~\ref{s6_highdim:sec:equal-its}--\ref{s6_highdim:sec:compute}).
On \Pone\ the lead of the PINN at the larger budget is of the size of the penalty bias of Deep Ritz
(Section~\ref{s6_highdim:sec:weights}). As in \citet{chen2020comparison}, neither loss dominates.
% src: data/s6_highdim_bound_numbers.json (n_networks_bound_valid; eta_d10_d20_all_networks; eta_by_d_all_networks; recheck_r2: pairs_ranked_like_error_equal_iterations 4/9); research_highdim_bound/results/decisions.json (H2: 21 of 57, max 5.456)
\item \emph{On the unit cube the error of a network is bounded by the two population terms of its
loss, but the bound is loose at small $d$.} The constant of the harmonic extension is of exact order $d^{-1/4}$
(Theorem~\ref{s6_highdim:thm:kappa}); the bound of Corollary~\ref{s6_highdim:cor:bound} held on all
133 networks evaluated, with efficiency 1.92 to 3.44 at $d=10$ and 20, but the pre-registered target
of an efficiency of at most 3 for $d\ge5$ failed (21 of 57 networks), and at $d=2$ the efficiency
reached 24.8. It is a certificate, not an error estimator. A bound of this form already follows from
Payne's classical constant; the theorem adds the order in $d$.
% src: research_highdim_remedies/results/decision.json; research_highdim_remedies/results/decision_p5.json; research_highdim_remedies/results/summary.json (P2/d20 lift3c, rep3; presolve and lift3c medians); research_highdim_remedies/results/summary_exploratory.json (P2/d20 lift1); research_highdim_remedies/results/attribution.json (P2 lift1/lift3c); research_highdim_remedies/results/summary_p5.json (P5/d20 presolve, lift3c); research_highdim_remedies/results/checks.json (C5_floors poisson/d20: exact_lift3 3.9e-3)
\item \emph{Neither cheap remedy is general at $d=20$.} By the pre-registered rule the Legendre
features (lift3c) were a general remedy on the four pre-registered problems, but they and the fixed
affine part were less accurate than plain on \Paltridge, which was added afterwards and was not a
blind test; the affine part also hurt on \Pridge. The features left the plateau of \Ptwo, whose
solution they contain up to $\sci{3.9}{-3}$, as did the centred input alone (exploratory, with
less gain) and, for Deep Ritz, the featureless control rep3 (with more); how much of each gain the
features explain differs between problems (Section~\ref{s6_highdim:sec:remedies}).
% src: data/closed_forms.json (floors); data/s4_lowdim_timeslices.json (summary: n_seeds_slice_t1_above_1 = 10 in all four cells)
\item \emph{How an error is reported changes the conclusion.} A relative $L^2$ error needs the error
of a trivial predictor beside it (Proposition~\ref{s2_problems:prop:floors}), and for a decaying
solution the space--time error hides the late slices: all 40 runs of Section~\ref{s4_lowdim:sec:timeslice}
are more than 100 per cent in error on the slice $t=1$, at a space--time error of about 0.1 per cent
after \AdamLBFGS.
\end{enumerate}

\subsection{A checklist for small PINN studies}\label{s8_discussion:sec:checklist}

\begin{enumerate}[label=\arabic*.,leftmargin=*]
\item Unit-test every residual operator on an exact solution and on a non-solution, and never
let a missing derivative default to zero (Section~\ref{s7_pitfalls:sec:autograd}).
\item Count the initial and boundary conditions and say why the solution is unique
(Section~\ref{s7_pitfalls:sec:oneface}).
\item Fix and report seeds, at least five, preferably ten, per configuration, and compare
paired by seed.
\item Report held-out relative errors, not training losses, per time slice for evolution
problems (Section~\ref{s4_lowdim:sec:timeslice}), with the error of a trivial predictor beside
them (Proposition~\ref{s2_problems:prop:floors}).
\item Run a classical baseline on the same evaluation set, after checking it against an
exact solution (Section~\ref{s5_hard:sec:fd}).
\end{enumerate}

\subsection{Limitations}\label{s8_discussion:sec:limitations}

\emph{Budgets.} The benchmark uses 256 to 1024 interior points and 3000 iterations. The
strategy effects belong to this scarce-point regime.

\emph{Statistics.} \PWone, \PWtwo\ and \PLthree\ have five seeds, for which the paired
test cannot go below $p=0.0625$, so ``no significant difference'' is weak evidence there.
The $p$-values are not corrected for multiple comparisons, and the rule by which we call
an effect significant (all ten seeds agree; Section~\ref{s3_methods:sec:stats}) was fixed
after the results were known.
% src: notes/VERIFICATION.md (sec. 2.5, caveat 2)

\emph{Verification.} The re-checks were made with separately written code within the same
project, on the same machine and library versions; they are not a replication by a third
party. The longer runs, the control experiments, the time-slice result, the 16\,000-iteration
runs at $d=20$ and the demonstration of Section~\ref{s7_pitfalls:sec:autograd} rest on the code of
the experiments alone (Section~\ref{s3_methods:sec:verif}).

\emph{Scope.} All problems are linear with constant coefficients on boxes, and their
solutions are smooth except at the discontinuous data of \PLthree. Each problem has one
tanh architecture, and Sections~\ref{s4_lowdim:sec} and~\ref{s5_hard:sec} use unit loss
weights. Uniqueness is proved in a stated class (classical solutions; for \PLthree, the
solutions that attain the datum in mean square); continuous dependence on the data is not
proved.

\emph{Unfinished problems.} \PWtwo\ and \PLthree\ are not converged at 3000 iterations.
The longer runs on \PWtwo\ (three seeds, one strategy) ended on the tolerance of the
optimiser in single precision, below one per cent but not on a plateau
(Section~\ref{s5_hard:sec:w2}). \PWone\ is not solved at the network size and budgets
used (Section~\ref{s5_hard:sec:w1}).
% src: research_benchmark/results/runs/wave2d_s0-5.json, laplace3d_s0-5.json (median held-out error of the 25 runs still falling between iterations 2750 and 3000: 5.81e-2 to 4.58e-2 and 4.56e-2 to 4.27e-2)
% src: data/s5_hard_wave2d_long.json (summary.per_seed: stop, adam_lbfgs_test_ratio_9500_over_8500); notes/VERIFICATION.md (sec. 2.5, caveat 6)

\emph{Implementation.} Most multi-seed runs used single precision on a GPU backend. A stack
repeated on the same machine reproduces its results bit for bit, but the result of one
network depends on the stack in which it is trained (its size and composition; the mechanism,
presumably rounding, was not investigated; Supplementary Section~\ref{sC_details:sec:repro}); other hardware was not tested. The \LBFGS\ of these runs
is a custom stacked implementation whose final errors were 4 to 10 per cent above those of the
reference optimiser, in geometric mean, in the validation of
Supplementary Section~\ref{sC_details:sec:validation}.
% src: research_benchmark/results/validate_batched.json (ratio stacked/single: geometric mean 1.10 on heat1d, 1.04 on laplace2d); notes/VERIFICATION.md (sec. 2.5, caveats 3 and 4)

\emph{Timings} come from three measurements on a loaded shared machine, with different
drivers, and are good to an order of magnitude only.
% src: data/checks/machine_load.json; notes/VERIFICATION.md (sec. 2.3, item 7; sec. 3.3, item 4)

\emph{Dimension study.} There are three seeds per cell, which understates the spread
(Table~\ref{s6_highdim:tab:vdim}) and allows no test, and no network was trained for
$d>20$. One architecture and one learning-rate schedule serve every $d$; the schedule
itself arrests the descent, so the errors at a fixed budget include its effect. The
weights were tuned at $d=5$ with one run per value, and $\lam$ and $\bet$ are normalised
differently in $d$ (Remark~\ref{s3_methods:rem:consistent}). Both solutions are simple in
structure, a sum of products of pairs of coordinates and a sum of functions of one
coordinate, and that of \Ptwo, with zero normal derivative, favours Deep Ritz. Boundary
conditions are imposed by penalty only. The inputs of the dimension study are not centred;
% src: research_highdim_remedies/results/summary_exploratory.json (P1, P2 at d = 20: paired plain/lift1 medians 0.92, 1.78, 1.57, 1.95; per seed 0.83 to 2.27); research_highdim_remedies/results/summary.json (plain)
centring alone (exploratory, Section~\ref{s6_highdim:sec:remedies}) changes the errors at $d=20$
and 4000 iterations by median factors of 0.92 to 1.95 (0.83 to 2.27 per seed), makes the PINN more
accurate than Deep Ritz on \Pone\ and takes both methods below the affine floor on \Ptwo. The 16\,000-iteration runs exist for $d=10$ on
\Pone\ and for $d=20$ only, and the equal-compute comparison for $d=10$ and, in the re-check,
for \Pone\ at $d=20$ with three seeds.
% src: notes/VERIFICATION.md (sec. 3.4; sec. 3.3, item 6)

\emph{Stability constant and remedies.} Theorem~\ref{s6_highdim:thm:kappa} and the bound of
Corollary~\ref{s6_highdim:cor:bound} are proved for the unit cube only, and whether $U_d$ is sharp
is open; the bound holds for the exact norms and was evaluated with Monte-Carlo estimates of them,
with margins of at least 75 standard errors on the 88 saved networks and 92 on the 21 networks of
the second round of the re-check; for the 24 networks of its first round (efficiency at least 1.97)
no standard errors were recorded. The remedies were tested with the network, optimiser,
schedule and weights of Section~\ref{s6_highdim:sec}, with five seeds at $d=20$ and three at $d=5$
and $10$, and with CPU-time budgets measured on a loaded shared machine. Both remedies, the centred
input and, for Deep Ritz on \Pone, rep3 had been tried on \Pone\ and \Ptwo\ in exploratory pilots
before the freeze, so their direction on these two problems was known and only \Pridge\ and
\Pcospair\ were held out; the pre-registered controls lift3u and rep3 were run for Deep Ritz only;
\Paltridge\ was not a blind test; and the second round of the re-check, which added \Psinlin\
and \Pexpcos, is descriptive, and three planned cells of it were not run (Supplementary
Section~\ref{S5:sec:remedies}).
% src: research_highdim_bound/results/decisions.json (H1: min_z_D hd 75.49); data/s6_highdim_bound_numbers.json (recheck_r2: min_z 92.57; n_networks_bound_valid: recheck_round1_eta_range from 1.974); research_highdim_bound/check/out/c4_fresh.jsonl (no standard-error fields); data/s6_highdim_remedies_recheck.json (missing_runs_of_plan; load1_range); notes/VERIFICATION.md (sec. 4.3, pilots before the freeze)

\subsection{What was not tried}\label{s8_discussion:sec:nottried}

Apart from the two cheap remedies of Section~\ref{s6_highdim:sec:remedies}, none of the
published remedies for the failures seen here was tried: adaptive loss
weighting \citep{wangteng2021,wangyu2022}, causal training \citep{wangsankaran2024},
Fourier-feature inputs \citep{wangwang2021}, or trial functions that satisfy the initial
and boundary conditions exactly \citep{lagaris1998,sheng2021,sukumar2022exact}; on the cube the
last would also remove the boundary term of the bound of Section~\ref{s6_highdim:sec:bound} and
the penalty bias. For Deep Ritz we tried neither continuation in $\bet$ nor a consistent boundary
treatment of Nitsche type~\citep{nitsche1971}, which exists in a neural
version~\citep{liao2021nitsche}. Non-linear and variable-coefficient equations, domains other than
boxes, and more seeds are the natural next steps.

\subsection{Conclusion}\label{s8_discussion:sec:conclusion}

We posed eight linear problems so that each has a unique solution in a stated class
(Proposition~\ref{s2_problems:prop:unique}), measured seeded held-out errors of PINNs and the Deep
Ritz method against exact solutions and, in low dimension, finite-difference baselines, and checked
the benchmark and the dimension study in part with separately written code. On these problems an \LBFGS\
phase mattered most; after it, and once the strip next to the initial line had residual points, the
collocation strategy changed the error by at most about a factor of two in geometric mean
($\Nr\ge256$); a classical solver was far faster; a unit-weight PINN of the size used did not solve
the two-mode wave problem; and whether the PINN or Deep Ritz is ahead in higher dimension depended
on the budget, on what is counted, on the input scaling and, on the Laplace problem, on the penalty
bias. On the unit cube a classical stability constant now has explicit two-sided bounds of the right
order in $d$; with it, as already with Payne's classical constant, the two population terms of the
PINN loss bound the $L^2$ error of any $C^2$ network, and the theorem sharpens that bound. The bound held on every network
evaluated but is loose. Of two cheap remedies at $d=20$, the fixed affine part did not remove the
plateau of the Poisson problem, and the Legendre features, which contain its solution up to 0.4 per
cent, did, as did centring the input alone (exploratory); neither remedy improved the accuracy on
every problem. Finally, the two pitfalls of
Section~\ref{s7_pitfalls:sec} show that a small training loss is not evidence of a correct solution;
each is caught by an entry of the checklist above.
